\chapter{Aplicación de consulta}
\label{ch:aplicacion}
% Fuentes: STREAMLIT_CODE_GUIDE.md; USER_GUIDE.md §§6--8;
% FINAL_STREAMLIT_PRODUCT_ALIGNMENT.md; APP_PRODUCT_SPEC.md.
\section{Objetivo, usuarios y arquitectura}
\pendiente{Definir Streamlit y su uso para crear aplicaciones en Python.
Describir usuarios potenciales, lectura exclusiva de
Gold validado y separación del proceso de extracción. No afirmar aceptación
pública de despliegue sin su evidencia de entrega.}
La guía versionada \path{docs/USER_GUIDE.md}, presentada en el
capítulo~\ref{ch:metodologia}, acompaña la aplicación: explica su instalación,
las vistas, los filtros, la ficha y las condiciones de las operaciones
controladas. La memoria explica el propósito del producto; el manual permite
seguir los pasos concretos para utilizarlo.
\section{Resumen, catálogo y filtros}
\pendiente{Presentar las vistas reales, los dos indicadores, filtros
compartidos, gráficos y mapas. Contar BOE distintos y proyectos por territorio;
los conteos territoriales no son sumables. Capturas pendientes de revisión.}
\pendiente{CAPTURAS C01 y C02 — Reservar vista general y mapa/filtros,
respectivamente. Mostrar contexto de consulta y cómo cambia la selección;
no duplicar dos vistas casi idénticas. Capturar después de la revisión de
publicación, con texto legible, caption y referencia desde esta sección.}
\section{Ficha del proyecto y cronología}
\pendiente{Describir territorio, publicaciones, actuaciones y acceso al BOE.
Explicar la lectura de la cronología y la preservación de incertidumbre.
Añadir una captura revisada del ejemplo del corpus.}
\pendiente{CAPTURA C03 — Ficha de proyecto y publicaciones: mostrar una
consulta verificable y su cronología, preferiblemente Don Rodrigo II de 2026.
Incluir caption y referencia desde el texto. Conservar entre dos y tres
capturas finales en total; no generar imágenes en esta integración.}
\section{Metodología, reporte y límites}
\label{sec:reporte-publico}
% streamlit_app.py::_render_report_channel; app_reporting.py; USER_GUIDE.md.
El usuario dispone del control «Reportar posible error» en la barra lateral,
debajo de la navegación. Cuando hay una dirección de revisión configurada,
el enlace utiliza \texttt{mailto:} para abrir un borrador en el cliente de
correo del usuario. Incluye la vista y los filtros activos; desde una ficha
añade contexto del proyecto y de su última publicación observada. El usuario
debe completar la descripción y enviar el mensaje desde su cliente de correo.

La aplicación no envía ni almacena el reporte: no crea una fila en una tabla
ni una entrada en la cola de revisión. Tampoco ofrece un backend de incidencias
ni integración automática con la Admin CLI. Sin una dirección válida,
el control queda deshabilitado y muestra que el canal no está configurado.
El destino se proporciona mediante configuración del operador; su disponibilidad
en una instalación debe comprobarse antes de presentarlo como canal operativo.

Si un mensaje recibido motiva una corrección, una persona debe trasladar el
caso al procedimiento técnico correspondiente, revisar la evidencia y preparar
una decisión trazable antes de regenerar los datos. Este paso es externo a
Streamlit y está sujeto a los límites de la sección~\ref{sec:revision-humana}.
La integración persistente de reportes con una cola unificada se propone
como trabajo futuro en la sección~\ref{sec:cola-unificada-futura}.

\pendiente{Completar la descripción del apartado metodológico. Distinguir esta
aplicación de consulta de la interfaz auxiliar de anotación del holdout,
descrita en el capítulo~\ref{ch:evaluacion}.}
